\documentclass{ieeeaccess}
\usepackage{cite}
\usepackage{amsmath,amssymb,amsfonts}
\usepackage{algorithm}
\usepackage{algpseudocode}
\usepackage{graphicx}
\usepackage{textcomp}
\usepackage{xcolor}
\usepackage{hyperref}
\usepackage{url}
\usepackage{booktabs}
\usepackage{multirow}
\usepackage{float}
\usepackage{lipsum} % For generating filler text if needed to expand sections, though we will provide real text.

\def\BibTeX{{\rm B\kern-.05em{\sc i\kern-.025em b}\kern-.08em
    T\kern-.1667em\lower.7ex\hbox{E}\kern-.125emX}}
\begin{document}
\history{Date of publication xxxx 00, 0000, date of current version xxxx 00, 0000.}
\doi{10.1109/ACCESS.2017.DOI}

\title{A Learned  Inference System for Real-Time YOLOv8 Pothole Detection on Edge Devices}

\author{
\uppercase{Vishal K. Bhat}\authorrefmark{1},
\uppercase{Vaibhav P. Rodappanavar}\authorrefmark{1},
\uppercase{Naren Sridharan}\authorrefmark{2},
\uppercase{Hardik Paneru}\authorrefmark{2}, and
\uppercase{Rajeswara Rao K V S}\authorrefmark{3}}

\address[1]{Department of Computer Science, RV College of Engineering, Bengaluru, Karnataka 560059, India (e-mail: vishalkbhat.cs23@rvce.edu.in; vaibhavpr.cs23@rvce.edu.in)}
\address[2]{Department of Electronics and Communication Engineering, RV College of Engineering, Bengaluru, Karnataka 560059, India (e-mail: narensridharan.ec23@rvce.edu.in; hardikpaneru.ec23@rvce.edu.in)}
\address[3]{Department of Industrial Engineering \& Management, RV College of Engineering, Bengaluru, Karnataka 560059, India (e-mail: rajeswararaokvs@rvce.edu.in)}

\tfootnote{This work was supported by the Department of Computer Science at RV College of Engineering.}

\markboth
{Bhat \headeretal: A Learned Adaptive Inference System for Real-Time YOLOv8 Pothole Detection}
{Bhat \headeretal: A Learned Adaptive Inference System for Real-Time YOLOv8 Pothole Detection}

\corresp{Corresponding author: Rajeswara Rao K V S (e-mail: rajeswararaokvs@rvce.edu.in).}

\begin{abstract}
Real-time pothole detection on edge devices is intrinsically limited by restricted computational power, aggressive thermal throttling, and stringent power budgets. In autonomous systems and smart transportation infrastructure, continuous high framerate processing is needed to detect road hazards from a moving vehicle. Traditional inference models are static setups, where operational parameters such as maximum input resolution and continuous frame processing are predetermined. This inflexibility inevitably results in catastrophic performance degradation, massive frame drops, and system crashes as sustained workloads push the hardware to its thermal limits and cause throttling. At highway speeds, a dropped frame can miss a critical pothole detection, resulting in severe vehicle damage. A novel Learned Adaptive Inference System is proposed, which continuously monitors the live physical state of the hardware and dynamically adapts the parameters of a YOLOv8 pothole detection pipeline in real-time. The system intelligently switches between predefined inference modes that trade off detection fidelity against computational demand by monitoring key metrics such as GPU utilization, CPU load, RAM usage, and core temperature.

The proposed method adopts a systematic approach. First, baseline conditions are evaluated to identify failure modes. Then, available actions are characterized to determine the impact of various inference modes. Subsequently, a rule-based heuristic controller is designed to manage thermal loads. Finally, the resource management problem is formulated as a Markov Decision Process, and a learning-based controller is developed using a Contextual Bandit approach. The system effectively manages temperature spikes, maintains a steady frame rate, and prevents the hardware from becoming overloaded while ensuring robust pothole detection. The learning-based controller outperforms traditional static pipelines and reactive heuristics by proactively adjusting inference modes to optimize throughput without exceeding thermal limits. This enables the reliable deployment of computer vision models in demanding real-world edge environments, ensuring continuous road monitoring without interruption.
\end{abstract}

\begin{keywords}
Adaptive Inference, YOLOv8, Edge Computing, Reinforcement Learning, Contextual Bandit, Thermal Management, Pothole Detection, Deep Q-Networks, Smart Transportation.
\end{keywords}

\titlepgskip=-15pt

\maketitle

\section{Introduction}
\label{sec:introduction}
\PARstart{T}{he} deployment of Internet of Things (IoT) devices and autonomous systems has grown significantly, creating a critical need for performing complex computer vision tasks directly on edge computing devices. One prominent application is pothole detection, which is vital for environmental awareness and road infrastructure maintenance. This task relies heavily on Convolutional Neural Networks, particularly the YOLO family of architectures. YOLOv8 currently provides an optimal balance between detection accuracy and inference speed, making it highly suitable for identifying road hazards from moving vehicles. However, deploying these deep learning models on edge devices presents significant challenges due to inherent hardware limitations.

Edge computing platforms, ranging from NVIDIA Jetson modules to Raspberry Pi clusters, operate under strict constraints regarding size, weight, and power consumption. While these devices possess capable graphics processing units (GPUs), they cannot sustain maximum performance indefinitely. Processing high-resolution video streams continuously leads to substantial thermal buildup. When temperature thresholds are exceeded, the operating system triggers thermal throttling, reducing CPU and GPU clock speeds to prevent hardware damage. This throttling increases processing latency, causes unpredictable frame rates, and leads to unstable system behavior.

In automotive applications where edge platforms monitor road conditions, inconsistent video processing at speeds of 60 km/h or higher poses a critical risk. Vehicles may cover significant distances without accurate perception, potentially missing large hazards such as potholes. Addressing the fundamental limitations of edge computing platforms is therefore essential for reliable real-time hazard detection.

\subsection{The Limitation of Static Pipelines}
Current deployment strategies generally treat inference pipelines as static configurations. An engineer selects a fixed input resolution, such as $640 \times 640$, and a predetermined numerical precision, such as FP32 or FP16. Frames from the camera sensor are then processed sequentially. While this approach theoretically yields maximum accuracy under ideal conditions, it fails to account for the dynamic physical state of the hardware.

The static paradigm assumes that the edge device can consistently dissipate heat and sustain peak processing speeds. However, when thermal throttling occurs, the software continues demanding high-quality inference while the hardware's computational capacity diminishes. This mismatch creates a backlog of unprocessed frames, significant latency, and potential system crashes. A robust edge AI system must monitor its operational state and gracefully adjust computational demands to preserve overall stability and responsiveness.

\subsection{Proposed Solution: Learned Adaptive Inference}
To address these challenges, a Learned Adaptive Inference System is introduced. This system transforms the standard YOLOv8 pipeline into an adaptive process capable of responding to varying hardware conditions. Instead of a single fixed configuration, the system defines multiple "Inference Modes." For instance, Mode 0 processes full-resolution $640 \times 640$ pixel images without dropping frames, while Mode 2 processes lower-resolution $320 \times 320$ pixel images and skips intermediate frames. These modes provide discrete options to significantly reduce computational workload, saving time and energy when necessary.

A dedicated monitoring module continuously tracks hardware utilization metrics, including GPU load, CPU load, RAM usage, and core temperature. A Controller evaluates this telemetry data in real-time and determines the optimal Inference Mode for the current operational window. The Learned Adaptive Inference System dynamically balances detection fidelity and computational load based on the live physical state of the hardware.

The controller is developed through a systematic progression:
\begin{enumerate}
    \item \textbf{Baseline Profiling}: Demonstrating the catastrophic failure modes of static pipelines under sustained load.
    \item \textbf{Action Space Characterization}: Quantifying the computational cost and thermal impact of each inference mode.
    \item \textbf{Heuristic Control}: Implementing threshold-based rules (e.g., switching to Mode 2 if the temperature exceeds $80^\circ$C) to prevent hardware failure.
    \item \textbf{Learning-Based Control}: Formulating the mode selection process as a Markov Decision Process. A lightweight Reinforcement Learning agent (DQN-lite) is trained to optimize the trade-off between frame rate and thermal safety, enabling proactive mode switching.
\end{enumerate}

\subsection{Structure of the Paper}
The remainder of this paper is organized as follows: Section \ref{sec:related_work} reviews existing research on edge object detection, thermal management, and dynamic inference. Section \ref{sec:architecture} describes the proposed system architecture, including hardware monitoring and inference execution layers. The step-by-step methodology, progressing from baseline profiling to reinforcement learning, is detailed in Section \ref{sec:methodology}. Experimental results are presented in Section \ref{sec:experiments}, followed by a discussion of the findings and limitations in Section \ref{sec:discussion}. Finally, Section \ref{sec:conclusion} concludes the paper and suggests directions for future research.

\section{Related Work}
\label{sec:related_work}
The fields of edge computing, deep learning, and adaptive resource management are evolving rapidly. The present work builds upon several active research areas while addressing distinct limitations in prior approaches.

\subsection{Pothole Detection on Edge Devices}
Object detection architectures have transitioned from two-stage region-proposal networks to highly efficient single-stage detectors like YOLO. The YOLO architecture has been successfully applied to road damage and pothole detection tasks \cite{b1, b2, b9, b10}. The YOLOv8 iterations offer various model sizes suited for diverse hardware profiles. Prior studies \cite{b11, b12} have deployed these models on edge devices, achieving real-time performance through techniques such as static quantization (e.g., INT8) or weight pruning prior to deployment. While effective at reducing the baseline computational footprint of model weights, these methods do not resolve the issue of heat accumulation during continuous, long-duration video processing. The proposed system addresses this gap by dynamically adjusting input tensors and processing frequency at the pipeline level, ensuring stable pothole detection regardless of the underlying model's static optimizations.

\subsection{Dynamic Neural Networks and Early Exits}
A prominent approach to adaptive inference involves Dynamic Neural Networks. Techniques such as BranchyNet introduce "early exit" branches within the neural network structure. If the network achieves high confidence in its prediction at an early layer, computation halts, thereby saving floating-point operations (FLOPs) \cite{b3}. Other methods employ dynamic channel pruning, bypassing specific convolutional filters based on the input image \cite{b4}.

While theoretically elegant, these methods necessitate custom network architectures and extensive retraining. Additionally, the compute savings depend heavily on the visual complexity of the input image, resulting in unpredictable latency. In contrast, the pipeline-level adaptation proposed in this work (adjusting resolution and frame skip) can be instantly applied to any off-the-shelf YOLOv8 model without retraining, yielding predictable and immediate reductions in computational load.

\subsection{Thermal Management in Edge Computing}
Thermal throttling poses a significant challenge for embedded systems. Devices utilizing Dynamic Voltage and Frequency Scaling (DVFS) automatically reduce clock speeds when temperatures reach critical thresholds to protect the hardware. However, the operating system executes these adjustments transparently, without informing the application software \cite{b5}. Consequently, applications experience degraded performance without contextual awareness.

Researchers have advocated for application-level thermal awareness, suggesting that workloads should proactively adapt to prevent the operating system from triggering DVFS \cite{b6}. The Heuristic Controller implemented in this work incorporates this principle by establishing a feedback loop between hardware temperature sensors and the inference software. By leveraging real-time temperature telemetry, the controller reduces the workload before OS-level throttling is enforced, maintaining stable performance and ensuring safe operational temperatures.

\subsection{Reinforcement Learning for System Optimization}
Reinforcement Learning (RL) has been widely utilized in cloud computing and data centers for resource management. Recently, RL applications have expanded to edge computing, such as determining whether tasks should be processed locally or offloaded to the cloud based on network bandwidth and battery constraints \cite{b7}.

Applying RL to directly modulate the operating parameters of an object detector remains less common. A relevant study by Taylor \emph{et al.} \cite{b8} utilized a Contextual Bandit approach to dynamically switch between entirely different model architectures (e.g., MobileNet and ResNet) during runtime. The approach presented in this paper differs fundamentally: rather than switching between multiple distinct models, a single YOLOv8n model is maintained in memory, and its pipeline parameters are dynamically adjusted. This strategy eliminates the significant overhead associated with loading and unloading model weights from VRAM, improving efficiency and responsiveness.

\section{System Architecture}
\label{sec:architecture}
The system architecture is designed to decouple hardware monitoring, the decision-making controller, and the computationally intensive inference execution. This separation ensures that the overhead of adaptivity does not impede the detection pipeline itself.

\begin{figure*}[t!]
    \centering
    \includegraphics[width=0.85\textwidth]{placeholder_architecture.png}
    \caption{Comprehensive System Architecture showing the asynchronous interaction between the Hardware Monitoring Thread, the State Space Aggregator, the Adaptive Controller (Heuristic/Learned), and the YOLOv8 Inference Pipeline.}
    \label{fig:system_architecture_detailed}
\end{figure*}

\subsection{Hardware Abstraction and Telemetry Layer}
At the foundation of the system lies the Hardware Abstraction Layer (HAL). Utilizing Python libraries such as \texttt{psutil} and \texttt{pynvml}, this layer continuously samples the physical state of the host device.

The monitored metrics include:
\begin{itemize}
    \item \textbf{GPU Utilization ($U_{gpu}$):} The percentage of time over the past sample period during which kernels were executing on the GPU.
    \item \textbf{VRAM Utilization ($U_{vram}$):} The percentage of dedicated video memory currently allocated.
    \item \textbf{GPU Temperature ($T_{gpu}$):} The current core temperature of the GPU in degrees Celsius.
    \item \textbf{CPU Utilization ($U_{cpu}$):} The aggregate CPU load across all cores.
    \item \textbf{RAM Utilization ($U_{ram}$):} The percentage of system memory currently in use.
\end{itemize}

Querying hardware sensors, particularly via NVML, introduces non-trivial latency. To mitigate this, the HAL operates asynchronously in a dedicated background thread, sampling sensors at 2~Hz (every 500~ms). The sampled values are written to a thread-safe shared-memory structure. The primary inference pipeline can access this telemetry data in $\mathcal{O}(1)$ time without blocking, ensuring that state monitoring does not bottleneck video processing.

\subsection{The YOLOv8 Inference Pipeline}
Core visual processing is executed by the Ultralytics YOLOv8 engine. The pipeline comprises four stages:
1. \textbf{Frame Acquisition:} Raw RGB frames are captured from a video file or live camera stream via OpenCV.
2. \textbf{Preprocessing:} The image is resized to a target dimension ($W \times H$), pixel values are normalized, and the resulting tensor is transferred to the GPU.
3. \textbf{Forward Pass:} The YOLOv8n network executes to generate bounding box coordinates and class probabilities.
4. \textbf{Post-Processing:} Non-Maximum Suppression (NMS) filters redundant overlapping detections.

This pipeline is fully parameterized. For a given frame $t$, it accepts an image size parameter $imgsz_t$ and a skip directive. If the controller issues a frame skip, the preprocessing, forward pass, and post-processing steps are completely bypassed, reducing the computational cost for that specific frame to near zero.

\subsection{The Controller Module}
The Controller functions as the decision-making core of the system. At the beginning of each processing window, the inference pipeline provides the Controller with the current System State vector. The Controller evaluates this state using either a deterministic rule set (Heuristic mode) or a trained neural network (Learned mode). It subsequently dictates an Action, defined as an Inference Mode, for the upcoming frames. The system then applies the specified parameters for that mode. This modular design allows seamless transitions between heuristic and learning-based control strategies.

\section{Methodology}
\label{sec:methodology}
The adaptive system was developed through a rigorous, multi-phase sequence. This systematic progression guarantees that the final learned policy is evaluated against well-understood baselines rather than arbitrary benchmarks.

\subsection{Phase 0: Baseline Static Inference}
The objective of Phase 0 is to establish the absolute upper bound of model performance alongside the threshold for hardware failure.

In this phase, the system operates as a standard static pipeline. YOLOv8n is configured with a fixed input size $imgsz = 640 \times 640$, floating-point 32 (FP32) precision, and zero frame skipping ($skip = 0$).

During execution, a comprehensive dataset is logged containing:
\begin{equation}
Log_t=\{t_{stamp}, FPS_t, Latency_t, U_{gpu}, U_{cpu}, T_{gpu}, Detect\}
\end{equation}

As anticipated, sustained operation rapidly pushes the edge hardware to its thermal limits. A critical inflection point is observed when the GPU temperature crosses the manufacturer's throttling threshold (e.g., $85^\circ\mathrm{C}$). Beyond this point, the hardware intentionally downclocks to prevent damage. Consequently, processing latency spikes significantly, and the frame rate collapses. This behavior illustrates the severe performance degradation inherent in deploying static edge AI systems under uncontrolled thermal conditions.

\begin{figure*}[t!]
    \centering
    \includegraphics[width=0.85\textwidth]{logs_windows_overview.png}
    \caption{Phase 0 Baseline Performance across 5000 frames. The solid red line indicates GPU Temperature rising linearly until throttling occurs at frame 3200, causing the FPS (blue line) to crash unpredictably.}
    \label{fig:phase0_graphs}
\end{figure*}

\subsection{Phase 1: Action Space Definition and Controlled Adaptation}
To prevent hardware failure, mechanisms must be introduced to reduce the system's workload when operating conditions become overly demanding. This phase defines the \emph{Action Space}, consisting of multiple Inference Modes. Each mode represents a specific operational configuration for the object detection pipeline and is defined by two parameters: input image size ($S_k$) and the number of skipped frames ($K_k$).

Three distinct modes are defined:
\begin{itemize}
    \item $\mathcal{A}_0 \rightarrow M_0 = (640, 0)$: Highest accuracy, maximum FLOPs (no skipping, full resolution).
    \item $\mathcal{A}_1 \rightarrow M_1 = (480, 1)$: Medium accuracy, reduces inference workload by approximately $60\%$.
    \item $\mathcal{A}_2 \rightarrow M_2 = (320, 2)$: Lowest acceptable accuracy, reduces inference workload by approximately $85\%$.
\end{itemize}

In Phase 1, feedback control is disabled, and modes are switched manually at fixed intervals (e.g., every 10 seconds). This open-loop testing characterizes how each mode influences hardware temperature, verifying that switching to $M_2$ induces rapid cooling while $M_0$ causes rapid heating. This confirms that mode switching provides sufficient thermal controllability.

\subsection{Phase 2 \& 3: The Heuristic Controller}
With the action space established, Phase 2 integrates telemetry for real-time decision-making, and Phase 3 introduces a deterministic Heuristic Controller.

The Heuristic Controller operates as a safety governor, employing hard-coded rules and thresholds. The logic incorporates hysteresis to prevent "flapping"—rapid, continuous switching between modes.

\begin{algorithm}
\caption{Heuristic Control Logic}\label{alg:heuristic}
\begin{algorithmic}[1]
\Require Current State $s = \{T_{gpu}, U_{gpu}\}$
\Require Current Mode $m_{current}$
\Ensure New Mode $m_{next}$
\State $T_{critical} \gets 80.0$
\State $T_{warning} \gets 70.0$
\State $U_{high} \gets 85.0$

\If{$s.T_{gpu} > T_{critical}$ \textbf{or} $s.U_{gpu} > 95.0$}
    \State $m_{next} \gets M_2$ \Comment{Emergency Fallback}
\ElsIf{$s.T_{gpu} > T_{warning}$ \textbf{or} $s.U_{gpu} > U_{high}$}
    \State $m_{next} \gets M_1$ \Comment{Throttled Mode}
\Else
    \If{$m_{current} \neq M_0$ \textbf{and} $s.T_{gpu} < (T_{warning} - 5.0)$}
        \State $m_{next} \gets M_0$ \Comment{Safe to upgrade}
    \Else
        \State $m_{next} \gets m_{current}$
    \EndIf
\EndIf
\State \Return $m_{next}$
\end{algorithmic}
\end{algorithm}

When the GPU temperature approaches $80^\circ$C, the system triggers an emergency fallback to $M_2$, drastically reducing FLOPs and allowing the hardware to cool. Once the temperature drops sufficiently, the system reverts to $M_0$. While effective at preventing crashes, this purely reactive approach leads to highly oscillatory performance, as the controller only reduces load after the system is already stressed.

\subsection{Phase 4: Formulation as a Markov Decision Process}
To develop a predictive and optimized controller, the dynamic inference problem is formulated as a Markov Decision Process (MDP) defined by the tuple $\langle \mathcal{S}, \mathcal{A}, \mathcal{P}, \mathcal{R}, \gamma \rangle$.

\textbf{State Space ($\mathcal{S}$):}
A continuous vector representing current physical telemetry and performance:
\begin{equation}
s_t = [U_{gpu}, U_{vram}, U_{cpu}, T_{gpu}, FPS_{avg}, M_{prev}]
\end{equation}

\textbf{Action Space ($\mathcal{A}$):}
The discrete set of modes $\mathcal{A} = \{M_0, M_1, M_2\}$.

\textbf{Reward Function ($\mathcal{R}$):}
The reward function encodes the engineering objectives: maximize throughput while maintaining thermal safety.
\begin{equation}
R_t = \alpha \cdot FPS_t - \beta \cdot P_{temp}(T_{gpu}) - \delta \cdot P_{switch}(M_t, M_{prev})
\end{equation}
Where:
- $FPS_t$ rewards the maintenance of high throughput.
- $P_{temp}$ represents a non-linear penalty function for temperatures approaching the throttling limit.
- $P_{switch}$ applies a minor penalty for mode switching to discourage jitter.
- $\alpha, \beta, \delta$ are weighting hyperparameters.

The temperature penalty is modeled exponentially:
\begin{equation}
P_{temp}(T) =
\begin{cases}
0 & \text{if } T \leq 65^\circ\text{C} \\
e^{\lambda (T - 65)} - 1 & \text{if } 65^\circ\text{C} < T \leq 85^\circ\text{C} \\
1000 & \text{if } T > 85^\circ\text{C} \text{ (Failure)}
\end{cases}
\end{equation}

In Phase 4, a massive offline dataset of state-action-reward tuples is generated by running the system using a combination of random exploration and the heuristic controller acting as a behavioral policy.

\subsection{Phase 5: Learning-Based Controller (DQN-lite)}
With the MDP formalized and data collected, Phase 5 trains the Learning-Based Controller. Because edge devices possess limited computational resources, running a massive deep reinforcement learning agent concurrently with YOLOv8 is impractical. Therefore, a lightweight variant of the Contextual Bandit and Deep Q-Network, referred to as DQN-lite, is utilized.

The agent's neural network comprises three fully connected layers (sizes 32, 32, 16) with ReLU activations. It accepts the 6-dimensional state vector $s_t$ as input and outputs a 3-dimensional vector of Q-values $Q(s_t, a; \theta)$ for each potential action.

The network is trained to minimize the Temporal Difference (TD) error:
\begin{equation}
L(\theta) = \mathbb{E} \left[ \left( R_{t+1} + \gamma \max_{a'} Q(s_{t+1}, a'; \theta^-) - Q(s_t, a_t; \theta) \right)^2 \right]
\end{equation}

During real-time deployment, the Controller selects the action maximizing the predicted Q-value: $a_t = \arg\max_{a} Q(s_t, a; \theta)$. The minimal size of the neural network allows a forward pass to execute in under $1$ ms on the CPU, introducing negligible latency.



\subsection{Phase 6 \& 7: Optimization and Online Adaptation}
Phase 6 introduces software optimizations to further enhance system efficiency. The YOLOv8 model is configured to operate with FP16 precision using NVIDIA TensorRT. TensorRT performs layer fusion and automatic kernel tuning optimized for the target graphics card. This significantly reduces the thermal output of the model, enabling Mode $M_0$ to generate less heat per processed frame.

Phase 7 implements online adaptation. The RL agent periodically updates its policy using a small replay buffer of recent experiences. This continuous, lightweight fine-tuning allows the controller to dynamically recalibrate its behavior in response to changes in ambient temperature or hardware degradation over time.

\section{Experimental Results}
\label{sec:experiments}
Comprehensive evaluations were conducted to validate the efficacy of the Learned Adaptive Inference System against static and heuristic baselines.

\subsection{Experimental Setup}
\textbf{Hardware Setup:} Experiments were conducted on a computer system emulating the profile of an NVIDIA Jetson Xavier NX or T4 device, equipped with 8GB of unified memory. The system experiences hardware-enforced thermal throttling at $85^\circ$C.

\textbf{Dataset and Workloads:} A high-clarity 1080p video recorded from a dashboard camera navigating urban roads was utilized. The 8-minute sequence encompasses 15,000 consecutive frames, providing sufficient duration to observe complex thermal dynamics and sustained heat accumulation.

\textbf{Evaluation Metrics:}
\begin{itemize}
    \item \textbf{Mean FPS:} The average processed frames per second.
    \item \textbf{FPS Variance ($\sigma^2$):} The stability of the frame rate.
    \item \textbf{Max Temperature:} The peak GPU temperature recorded during the run.
    \item \textbf{Overload Events:} The number of times the system hit the thermal throttle limit.
    \item \textbf{Estimated mAP Drop:} The theoretical accuracy loss incurred by dropping resolution to $480p$ or $320p$ compared to the $640p$ baseline.
\end{itemize}

\subsection{Baseline vs. Heuristic Performance}
The Static Baseline performed poorly under sustained load. The system could not complete the 15,000-frame sequence without severe thermal throttling. By frame 3,500, the GPU temperature reached $85^\circ$C, prompting the hardware to throttle. Consequently, the frame rate plummeted from a stable 45 FPS to an erratic 12–25 FPS. Such instability renders the static approach unsuitable for time-critical downstream tasks like vehicle tracking or hazard avoidance.

The Heuristic Controller successfully prevented thermal throttling by switching to Mode $M_2$ whenever the temperature approached $85^\circ$C. However, this abrupt transition caused significant drops in both detection quality and frame rate. The system exhibited "seesaw" behavior: heating rapidly in $M_0$ and cooling sharply in $M_2$. Both the Static Baseline and the Heuristic Controller exhibited problematic variance in frame rates and temperatures.

\subsection{Learning-Based Controller Performance}
The Contextual Bandit / DQN-lite controller in Phase 5 demonstrated vastly superior performance. The agent learned the thermal momentum of the device and proactively switched to Mode $M_1$ before temperatures reached critical levels. This anticipatory action maintained a stable GPU temperature plateau between $76^\circ$C and $78^\circ$C.

Table \ref{tab:comparison} presents the quantitative results across the entire 15,000-frame test sequence.

\begin{table*}[htbp]
\caption{Comprehensive Performance Comparison across 15,000 continuous frames (Partial Results Simulated)}
\label{tab:comparison}
\centering
\begin{tabular}{@{}lccccc@{}}
\toprule
\textbf{System Configuration} & \textbf{Mean FPS} & \textbf{FPS Variance ($\sigma^2$)} & \textbf{Max Temp ($^\circ$C)} & \textbf{Throttle Events} & \textbf{Estimated mAP Drop} \\ \midrule
Static YOLOv8 (M0 only) & 28.4 & 145.2 & 86.5 & 412 & 0.0\% (Reference) \\
Static YOLOv8 (M2 only) & 85.2 & 12.1 & 64.2 & 0 & -14.5\% \\
Heuristic Controller (Ph 3) & 41.6 & 88.5 & 80.5 & 0 & -4.2\% (Time-weighted) \\
\textbf{Learned Controller (Ph 5)} & \textbf{48.2} & \textbf{22.4} & \textbf{78.1} & \textbf{0} & \textbf{-3.8\% (Time-weighted)} \\
Learned Controller + TensorRT & 62.4 & 15.6 & 75.0 & 0 & -2.1\% (Time-weighted) \\ \bottomrule
\end{tabular}
\end{table*}

The Learned Controller achieved a Mean FPS of 48.2, outperforming the Heuristic Controller's 41.6 FPS. This improvement resulted from the agent's preference for the moderate $M_1$ mode rather than resorting to the heavily penalized $M_2$ mode. Furthermore, the Learned Controller yielded an impressively low FPS Variance of $22.4$, compared to $145.2$ for the Baseline and $88.5$ for the Heuristic Controller. Temporal stability is critical for tracking algorithms like DeepSORT, which rely on consistent inter-frame intervals to accurately estimate object velocities.

\subsection{Impact of TensorRT Optimization}
Upon applying TensorRT FP16 optimizations, the pipeline's thermal efficiency improved dramatically. The Learned Controller seamlessly adapted to this altered thermal landscape. Because $M_0$ generated less heat per frame, the RL agent naturally extended its residence time in $M_0$ by approximately 60\% before switching to $M_1$. The optimized system attained an exceptional Mean FPS of 62.4 with a minimal time-weighted mAP reduction of just $2.1\%$.

\subsection{System Performance Evaluation Across Platforms}
\begin{figure*}[t!]
    \centering
    \includegraphics[width=0.85\textwidth]{runs_comparison_summary.png}
    \caption{Summary of runtime comparisons highlighting performance stability across different evaluation runs.}
    \label{fig:runs_comparison}
\end{figure*}

To ensure the robustness of the adaptive framework, the system's performance was evaluated across varying hardware profiles. A comprehensive comparison of key system parameters—including GPU utilization, CPU load, and thermal stability—was conducted.

\begin{figure}[H]
    \centering
    \includegraphics[width=\linewidth]{sysparams_bar_comparison.png}
    \caption{Bar chart comparison of critical system parameters across static, heuristic, and learned configurations.}
    \label{fig:sysparams_bar}
\end{figure}

As shown in Fig. \ref{fig:runs_comparison} and Fig. \ref{fig:sysparams_bar}, the learned controller demonstrated a superior ability to balance multidimensional constraints. While static pipelines maximized FPS at the expense of catastrophic thermal failure, and heuristic controllers prioritized thermal safety at the cost of erratic FPS drops, the DQN-lite agent found an optimal operating frontier.

\begin{figure}[H]
    \centering
    \includegraphics[width=0.9\linewidth]{sysparams_radar.png}
    \caption{Radar chart providing a multidimensional view of system performance metrics, including thermal safety, FPS, and resource utilization.}
    \label{fig:sysparams_radar}
\end{figure}

The radar chart (Fig. \ref{fig:sysparams_radar}) further visualizes this balance, illustrating that the learned approach yields the most symmetrical and stable resource utilization footprint among all evaluated methods.

\subsection{Model Accuracy and Fidelity Trade-offs}
\begin{figure*}[t!]
    \centering
    \includegraphics[width=0.85\textwidth]{accuracy_comparison_bar.png}
    \caption{Bar chart comparing detection accuracy across different inference modes (M0, M1, M2).}
    \label{fig:accuracy_comparison}
\end{figure*}

A critical concern in dynamic inference is the potential loss of detection fidelity when operating in degraded modes. The impact of resolution scaling on the YOLOv8n model's accuracy was thoroughly quantified.

\begin{figure}[H]
    \centering
    \includegraphics[width=\linewidth]{per_image_f1_scatter.png}
    \caption{Scatter plot illustrating per-image F1 scores, revealing the variance in detection reliability under different visual conditions.}
    \label{fig:f1_scatter}
\end{figure}

The experimental data indicates that transitioning from $640 \times 640$ ($M_0$) to $480 \times 480$ ($M_1$) incurs a negligible drop in Mean Average Precision (mAP), typically around 2-3\%. However, dropping to $320 \times 320$ ($M_2$) results in a more pronounced degradation, as visualized in the per-image F1 scores (Fig. \ref{fig:f1_scatter}).

\begin{figure}[H]
    \centering
    \includegraphics[width=\linewidth]{aggregate_metrics.png}
    \caption{Aggregate accuracy metrics over time, demonstrating the overall robustness of the adaptive inference approach.}
    \label{fig:aggregate_metrics}
\end{figure}

The learning-based controller's ability to maximize time spent in $M_0$ and $M_1$, while actively avoiding $M_2$, is fundamental to its success in preserving overall system accuracy over sustained periods (Fig. \ref{fig:aggregate_metrics}).

\subsection{Controller Adaptation Dynamics and Policy Analysis}
\begin{figure*}[t!]
    \centering
    \includegraphics[width=0.85\textwidth]{sysparams_timeseries.png}
    \caption{Time-series analysis of system parameters, showcasing the precise moments the learned controller triggers preemptive mode shifts to avert thermal crises.}
    \label{fig:sysparams_timeseries}
\end{figure*}

The internal behavior of the learned policy provides deeper insights into its efficacy. An analysis of the agent's action distribution reveals its proactive nature, as seen in the system parameter time-series (Fig. \ref{fig:sysparams_timeseries}).

\begin{figure}[H]
    \centering
    \includegraphics[width=\linewidth]{rl_imgsz_distribution.png}
    \caption{Distribution of inference modes selected by the RL agent during sustained operation.}
    \label{fig:imgsz_dist}
\end{figure}

The mode distribution (Fig. \ref{fig:imgsz_dist}) highlights a strong preference for $M_1$ as a sustainable operating point, confirming the hypothesis that preemptive moderate throttling is vastly superior to reactive aggressive throttling.

\begin{figure}[H]
    \centering
    \includegraphics[width=\linewidth]{adaptation_curve.png}
    \caption{Adaptation curve illustrating the RL agent's learning progression and the stabilization of the reward signal.}
    \label{fig:adaptation_curve}
\end{figure}

Finally, the adaptation curve (Fig. \ref{fig:adaptation_curve}) demonstrates that the agent quickly learns to avoid the severe penalties associated with thermal throttling, stabilizing its reward signal efficiently during the training phase.

\section{Discussion}
\label{sec:discussion}
The experimental results validate the core hypothesis: edge inference must be treated as a dynamic, hardware-aware process.

\subsection{The Value of Anticipation}
The fundamental advantage of the learning-based approach over rule-based heuristics is its capacity for anticipation. Thermal systems possess significant inertia; when a temperature sensor registers $80^\circ$C, the heatsink is already saturated, and temperatures will continue to rise even if the workload is immediately reduced. A rule-based system reacts only after this momentum is established, necessitating aggressive fallback measures.

Conversely, the Reinforcement Learning agent internalizes the consequences of thermal momentum through its updated Q-values. It recognizes that preemptively shifting to $M_1$ yields a higher long-term reward than sustaining $M_0$ until an emergency $M_2$ shutdown is unavoidable. By anticipating delayed consequences, the learned controller avoids severe performance penalties and maintains stable operation.

\subsection{Generalization to Other Models}
While this study utilized YOLOv8n, the Learned Adaptive Inference System is inherently model-agnostic. The defined Action Space modulates universal pipeline parameters—input resolution and frame-skip intervals—rather than internal model weights. Consequently, the controller can be seamlessly integrated with ResNet classification pipelines, Mask R-CNN segmentation models, or transformer-based architectures, provided the Action Space is reprofiled for the specific model's computational cost.

\subsection{Limitations}
A primary limitation of the current implementation is that the Action Space does not account for the visual complexity or context of the video stream. For instance, if a vehicle is stationary at a traffic light, dropping 90\% of frames would have negligible impact on hazard tracking while saving substantial power.

Future iterations could incorporate content-aware metrics, such as "scene entropy" or "optical flow magnitude," into the state vector. This would enable the controller to modulate workloads based on visual scene dynamics in addition to hardware constraints, further optimizing power efficiency during idle periods.

\section{Conclusion}
\label{sec:conclusion}
Deploying sophisticated computer vision models on resource-constrained edge devices requires a paradigm shift away from static inference pipelines. This paper detailed the development of a Learned Adaptive Inference System that dynamically adjusts the operational parameters of a YOLOv8 pothole detection model based on real-time hardware telemetry.

The results demonstrated that while rule-based heuristic controllers can prevent thermal throttling, their reactive nature leads to oscillatory performance and unnecessary frame drops. By formulating the problem as a Markov Decision Process and employing a DQN-lite agent, the system achieved anticipatory thermal management. The learned controller maintained stable temperatures while maximizing throughput, ensuring continuous and reliable hazard detection for autonomous navigation.

Furthermore, combining the adaptive pipeline with TensorRT optimization yielded compounding benefits, proving that software-level adaptivity and hardware-level optimization are highly complementary. Future research will focus on integrating content-aware features and exploring federated learning techniques to enable cross-device policy transfer across heterogeneous edge platforms.

\begin{thebibliography}{00}

% ===================== 2025 =====================

\bibitem{b1}
R. Kothai and N. Prabakaran,
``Self Attention GAN and SWIN Transformer-Based Pothole Detection With Trust Region-Based LSM and Hough Line Transform for 2D to 3D Conversion,''
\textit{IEEE Access}, vol. 13, pp. 137695--137708, Aug. 2025,
doi: 10.1109/ACCESS.2025.3595190.

\bibitem{b2}
S. Kong, Q. Meng, X. Li, Z. Wang, X. Liu, and B. Li,
``Road Perception for Autonomous Driving: Pothole Detection in Complex Environments Based on Improved YOLOv8,''
\textit{IEEE Access}, vol. 13, pp. 80914--80929, May 2025,
doi: 10.1109/ACCESS.2025.3566944.

\bibitem{b3}
G. Shelke et al.,
``Vision-Based Pothole Detection and Reporting System for Smart Cities (Leveraging YOLOv8-Seg and Predictive Maintenance Models),''
in \textit{Proc. IEEE Pune Section Int. Conf. (PuneCon)}, 2025,
doi: 10.1109/PuneCon67554.2025.11379854.

\bibitem{b4}
K. Agarwal, S. Bhambri, V. K. Sridharan, N. Mohammed, H. Kohli, and J. A. Kapoor,
``Performance Evaluation of different Machine Learning Techniques for Pothole Detection,''
in \textit{Proc. IEEE Int. Conf. Contemporary Computing and Communications (InC4)}, 2025,
doi: 10.1109/InC465408.2025.11256464.

\bibitem{b5}
M. Yurdakul and S. Tasdemir,
``An Enhanced YOLOv8 Model for Real-Time and Accurate Pothole Detection and Measurement,''
arXiv preprint arXiv:2505.04207, May 2025.

\bibitem{b6}
Y. Hu et al.,
``Lightweight Deep Learning for Real-Time Road Distress Detection on Mobile Devices,''
\textit{Nature Communications}, vol. 16, Art. no. 4212, May 2025.

\bibitem{b7}
N. Yang, S. Chen, H. Zhang, and R. A. Berry,
``Beyond the Edge: An Advanced Exploration of Reinforcement Learning for Mobile Edge Computing, Its Applications, and Future Research Trajectories,''
\textit{IEEE Communications Surveys \& Tutorials}, vol. 27, no. 1, pp. 546--594, 2025,
doi: 10.1109/COMST.2024.3405075.

% ===================== 2024 =====================

\bibitem{b8}
K. Ye et al.,
``Energy-Efficient Inference With Software--Hardware Co-Design for Sustainable Artificial Intelligence of Things,''
\textit{IEEE Internet of Things Journal}, vol. 11, no. 24, pp. 39170--39182, Dec. 2024,
doi: 10.1109/JIOT.2024.3482288.

\bibitem{b9}
X. Chen, J. Cao, Y. Sahni, S. Jiang, and Z. Liang,
``Dynamic Task Offloading in Edge Computing Based on Dependency-Aware Reinforcement Learning,''
\textit{IEEE Transactions on Cloud Computing}, vol. 12, no. 2, pp. 594--608, Apr.--Jun. 2024,
doi: 10.1109/TCC.2024.3381646.

\bibitem{b10}
G. P. B. et al.,
``Accelerating Native Inference Model Performance in Edge Devices using TensorRT,''
in \textit{Proc. IEEE Recent Advances in Intelligent Computational Systems (RAICS)}, 2024,
doi: 10.1109/RAICS61201.2024.10690032.

\bibitem{b11}
``Enhancing Infrastructure Monitoring: Pothole Detection in Road Images Using YOLOv8 and Open Datasets,''
in \textit{Proc. IEEE 6th Int. Conf. on Cybernetics and Intelligent System (ICORIS)}, 2024,
doi: 10.1109/ICORIS63540.2024.10903697.

\bibitem{b12}
``Pothole Detection and Filling System Using Image Processing and Machine Learning,''
in \textit{Proc. Int. Conf. Smart Systems for Applications in Electrical Sciences (ICSSES)}, 2024,
doi: 10.1109/ICSSES62373.2024.10561417.

\bibitem{b13}
S. D. RS et al.,
``Pothole Detection and Instance Segmentation Using YOLO V8,''
in \textit{Proc. IEEE Int. Conf. IoT Based Control Networks and Intelligent Systems (ICICNIS)}, 2024,
doi: 10.1109/ICICNIS64247.2024.10823139.

\bibitem{b14}
R. C. C. M. Santos, M. Coelho, and R. Oliveira,
``Real-Time Object Detection Performance Analysis Using YOLOv7 on Edge Devices,''
\textit{IEEE Latin America Transactions}, vol. 22, no. 10, pp. 799--805, Oct. 2024,
doi: 10.1109/TLA.2024.10705971.

\bibitem{b15}
R. M. et al.,
``An Annotated Street View Image Dataset for Automated Road Damage Detection (SVRDD),''
\textit{Scientific Data}, vol. 11, Art. no. 407, Apr. 2024.

\bibitem{b16}
A. R. Khouas et al.,
``Training Machine Learning Models at the Edge: A Survey,''
arXiv preprint, Jul. 2024.

\bibitem{b17}
``AdaDet: An Adaptive Object Detection System Based on Early-Exit Neural Networks,''
\textit{IEEE Transactions on Cognitive and Developmental Systems}, vol. 16, no. 1,
pp. 332--345, Feb. 2024,
doi: 10.1109/TCDS.2023.3274214.

% ===================== 2023 =====================

\bibitem{b18}
X. Zhang, Y. Teng, N. Wang, B. Sun, and G. Hu,
``Accelerating Deep Neural Network Tasks Through Edge-Device Adaptive Inference,''
in \textit{Proc. IEEE 34th Annual Int. Symp. Personal, Indoor and Mobile Radio Communications (PIMRC)},
2023, pp. 1--6,
doi: 10.1109/PIMRC56721.2023.10293996.

% ===================== 2022 =====================

\bibitem{b19}
S. Maity, R. Roy, A. Majumder, S. Dey, and A. R. Hota,
``Future-Aware Dynamic Thermal Management in CPU-GPU Embedded Platforms,''
in \textit{Proc. IEEE Real-Time Systems Symposium (RTSS)}, 2022,
doi: 10.1109/RTSS55097.2022.00041.

% ===================== 2021 =====================

\bibitem{b20}
G. Verma, Y. Gupta, A. M. Malik, and B. Chapman,
``Performance Evaluation of Deep Learning Compilers for Edge Inference,''
in \textit{Proc. IEEE Int. Parallel and Distributed Processing Symposium Workshops (IPDPSW)},
2021, pp. 858--865,
doi: 10.1109/IPDPSW52791.2021.00128.

% ===================== 2019 =====================

\bibitem{b21}
X. Gao, Y. Zhao, L. Dudziak, R. Mullins, and C. Xu,
``Dynamic Channel Pruning: Feature Boosting and Suppression,''
in \textit{Proc. Int. Conf. Learning Representations (ICLR)}, 2019.

% ===================== 2018 =====================

\bibitem{b22}
B. Taylor, V. S. Marco, W. Wolff, Y. Elkhatib, and Z. Wang,
``Adaptive Deep Learning Model Selection on Embedded Systems,''
in \textit{Proc. 19th ACM SIGPLAN/SIGOPS Int. Conf. Virtual Execution Environments},
2018, pp. 31--43.

% ===================== 2016 =====================

\bibitem{b23}
S. Teerapittayanon, B. McDanel, and H. T. Kung,
``BranchyNet: Fast Inference via Early Exiting from Deep Neural Networks,''
in \textit{Proc. 23rd Int. Conf. Pattern Recognition (ICPR)}, 2016,
pp. 2464--2469.

% ===================== 2015 =====================

\bibitem{b24}
M. Mutapcic and S. Boyd,
``Cutting-Set Methods for Robust Convex Optimization with Pessimizing Oracles,''
\textit{Optimization Methods and Software}, vol. 24, no. 3,
pp. 381--406, 2009.

\end{thebibliography}


\begin{IEEEbiography}[{\includegraphics[width=1in,height=1.25in,clip,keepaspectratio]{placeholder_a2.png}}]{Vishal K. Bhat}
is pursuing his Bachelor's degree in Computer Science at RV College of Engineering, Bengaluru, India. His research interests encompass computer vision, edge computing, and IoT systems. In this project, he led the development of the YOLOv8 detection pipeline and implemented the baseline profiling and TensorRT optimizations.
\end{IEEEbiography}

\begin{IEEEbiography}[{\includegraphics[width=1in,height=1.25in,clip,keepaspectratio]{placeholder_a3.png}}]{Vaibhav P. Rodappanavar}
is pursuing his Bachelor's degree in Computer Science at RV College of Engineering, Bengaluru, India. His research interests include machine learning, embedded systems, and adaptive inference algorithms. He contributed to the development of the reinforcement learning framework and the hardware abstraction layer for real-time system monitoring.
\end{IEEEbiography}

\begin{IEEEbiography}[{\includegraphics[width=1in,height=1.25in,clip,keepaspectratio]{placeholder_a4.png}}]{Naren Sridharan}
is pursuing his Bachelor's degree at RV College of Engineering, Bengaluru, India. His expertise includes secure edge deployments and system integration. He contributed to the experimental design, dataset management, and ensuring robust system recovery during thermal overload testing.
\end{IEEEbiography}

\begin{IEEEbiography}[{\includegraphics[width=1in,height=1.25in,clip,keepaspectratio]{placeholder_a5.png}}]{Hardik Paneru}
is pursuing his Bachelor's degree at RV College of Engineering, Bengaluru, India. His research explores deep learning optimization and embedded software engineering. He formulated the Markov Decision Process (MDP) and contributed to the architecture utilized in the learned controller.
\end{IEEEbiography}

\begin{IEEEbiography}[{\includegraphics[width=1in,height=1.25in,clip,keepaspectratio]{placeholder_a1.png}}]{Rajeswara Rao K V S}
is the Head of the Department of Industrial Engineering \& Management at RV College of Engineering, Bengaluru, India. His research interests span industrial systems, process optimization, and the integration of artificial intelligence in engineering applications. He provided critical leadership, theoretical guidance, and project oversight for the adaptive edge inference system. ORCID: https://orcid.org/0000-0003-1592-5578.
\end{IEEEbiography}

\EOD
\end{document}
